\documentclass[12pt,letterpaper]{article}

\usepackage[utf8]{inputenc}
\usepackage[T1]{fontenc}
\usepackage[english]{babel}
\usepackage{graphicx, geometry, xcolor}

\geometry{margin=2cm, top=2cm, bottom=2cm}

\begin{document}

\begin{titlepage}
  \thispagestyle{empty}
  %--------------- COMANDO LINEA----------------->
  \newcommand{\linea}{\rule{\linewidth}{0.5mm}}                 
  \center
  %<<<<<<<<<<<<<<<<<<<<<<<<<<<<<<<<<<<<<<<<<<<

  %Add logos
  \includegraphics[width=0.24\textwidth]{~/unam_logo.png} \hspace{1.5cm}
  \includegraphics[width=0.25\textwidth]{~/fc_logo.png} \\[1cm]
  \textbf{\Large UNIVERSIDAD NACIONAL AUTÓNOMA}\\[3mm]
  \textbf{\Large DE MÉXICO} \\[6mm]
  \textit{\Large Facultad de Ciencias}

  \vfill

  %---------------TÍTULO----------------->
  \linea
  \vfill \vspace{1cm}
  \textbf{\Large Programación de Dispositivos Móviles}\\[1.1cm]
  \textbf{\Large Cibersabiduría}\\[0.9cm]
  \linea \\
  \vfill
  \vspace{0.7cm}
  %Title of the Research
  \textbf{\Large  Práctica 3}\\[6mm]
  %<<<<<<<<<<<<<<<<<<<<<<<<<<<<<<<<<<<<<<<<<<<

  %Team
  \vfill
  \textit{\small Presenta:}\\
  \textbf{\large \textbf{Rojo Peña Manuel Ianluck}}\\[3mm]

  %Profesor y Grupo
  \textit{\small Profesor:}\\
  \textbf{Gustavo Arturo Márquez Flores}\\[3mm]
  \textit{\small Ayudante laboratorio:}\\
  \textbf{Jesús Iván Saavedra Martínez}\\[3mm]
  \small Grupo: 7138, 2026-2\\[3mm]
  \vfill

  %Fecha or Date
  \textit{\small Fecha de entrega:}\\
         {\large 13 de abril, 2026}
         
\end{titlepage}


\newpage

De manera formal, la cibersabiduría se define como:

\begin{quote}
  \textit{La capacidad de utilizar las tecnologías digitales de manera adecuada, haciendo lo correcto en el momento oportuno. La cibersabiduría es una meta-virtud basada en el concepto neoaristotélico de phronesis (frónesis) o sabiduría práctica, que permite a la persona comprender y actuar correctamente en sus circunstancias cotidianas.}
\end{quote}
  
  Como ejemplos prácticos podemos mencionar, en primer lugar, el uso correcto de internet, una herramienta que transformó al mundo al permitirnos comunicarnos globalmente, acceder a información de manera instantánea y traspasar esas barreras geográficas. Sin embargo, desde su popularización también hemos descubierto riesgos significativos, como el robo de identidad, las estafas en línea, la desinformación y la exposición de datos personales a terceros malintencionados. Esta dualidad nos ha enseñado que debemos ser cautelosos al navegar y compartir información, pues lo que puede ser una ventana al conocimiento también puede convertirse en una puerta de entrada a vulnerabilidades si no se usa con criterio. En segundo lugar, y de manera más reciente, está el uso de las inteligencias artificiales, donde mucha gente repite los mismos errores que se cometían en los inicios de internet, como entregar datos personales sin considerar quién los maneja ni con qué fines. Del mismo modo, muchos creen que la IA es una herramienta mágica que resolverá cualquier problema, cuando en realidad se necesitan conocimientos previos para saber qué pedirle y cómo pedírselo; no basta con hacer una consulta, sino que es necesario comprender el contexto y evaluar críticamente las respuestas que nos brinda, ya que confiar ciegamente en ella puede llevarnos a tomar decisiones basadas en información errónea o incompleta.\\

En lo personal, como ya mencioné, las inteligencias artificiales no son una herramienta mágica, más bien son un recurso que puede ofrecernos otra perspectiva, ayudarnos a centralizar ideas, generar una previsualización de un objetivo y facilitar que otros comprendan lo que una persona imagina. Sin embargo, al aparentar ser una máquina con respuesta para todo, mucha gente no las utiliza como se debe y terminan delegando completamente el trabajo a estas, tanto en el ámbito académico como en entornos laborales reales. Esto no es beneficioso, ya que uno debe comprender lo que la IA le está dando como respuesta para poder determinar si es útil o no. No podemos simplemente decir "\textit{haz este código}" y conformarnos; si le pedimos algo, debemos ser claros al momento de solicitarlo y rigurosos al analizar la respuesta, pues de lo contrario podríamos estar obteniendo resultados incorrectos sin siquiera saberlo.\\

El uso de asistentes de código me parece una herramienta bastante útil que puede agilizar el trabajo. Podemos emplearlos como un buen recurso de aprendizaje siempre y cuando verifiquemos que las respuestas sean válidas, ya que, por la forma en que estas inteligencias artificiales están programadas, pueden ser susceptibles a errores, aunque herramientas como \texttt{Claude} han minimizado ampliamente este margen de error en el campo de la programación. Son muy útiles porque nos brindan una perspectiva sobre nuestras ideas, sobre lo que queremos implementar, y nos permiten agilizar el proceso de explorar qué sucede si probamos cierta implementación. No obstante, esto también trae consigo un problema, algunos jefes consideran que pueden asignarnos más carga de trabajo al ver que con las IAs avanzamos más rápido, lo cual resulta irónico, ya que se supone que las usamos para agilizar nuestro desempeño, no para acumular más tareas.\\

\newpage

\textbf{Referecias:}\\

\begin{itemize}
\item Gómez-Gutiérrez, J. L., Fernández-Espinosa, V., \& Harrison, T. (2024). Un enfoque de cibersabiduría para la educación en ciudadanía digital. Percepciones de adolescentes españoles. Bordón Revista de Pedagogía, 76(2), 173-196.\\ https://doi.org/10.13042/bordon.2024.100155
\end{itemize}


\end{document}
